%======================================================================
\section{Valuation}\label{sec:valuation}
%======================================================================

\subsection{The pricing equation}

Let $r$ be the flat, continuously compounded risk-free rate used for discounting, $F$ the
face value, $f$ the coupon frequency and $t_1<t_2<\dots$ the coupon dates. The scheduled
coupon rate is $c(t_i)=c_0$ up to and including the first reset $T_1$ and
$c(t_i)=y^{\mathrm{ref}}+m$ after it, where $y^{\mathrm{ref}}$ is the five-year reference
rate. A coupon is paid if the bond has not been converted, called or written down, scaled
by the MDA fraction. The principal is repaid at the call, converted at the trigger, or
lost at PONV. The dirty price at time 0 is
\begin{equation}\label{eq:price}
\begin{split}
  V_0=\E\Biggl[\,&\sum_{t_i\le H}\1_{\{t_i<\tau_m,\ t_i\le\tau_c,\ t_i<\tau_P\}}
  \,\phi(C_{t_i},t_i)\,\frac{c(t_i)F}{f}\,e^{-rt_i}
  + \1_{\{\tau_c<\tau_m\wedge\tau_P,\ \tau_c\le H\}}\,F\,e^{-r\tau_c}\\
  &+ \1_{\{\tau_m<\tau_c\wedge\tau_P,\ \tau_m\le H\}}\,R_{\mathrm{conv}}F\,e^{-r\tau_m}
  + \1_{\{\tau_c\wedge\tau_m\wedge\tau_P>H\}}\,T_H\,e^{-rH}\Biggr],
\end{split}
\end{equation}
where $H$ is the simulation horizon and $T_H$ the value at $H$ of a bond still alive,
\begin{equation}\label{eq:tail}
  T_H = \frac{(y^{\mathrm{ref}}+m)F}{r+\lambda\, s_H^+/s_0}.
\end{equation}
The tail is the value of a perpetuity paying the reset coupon continuously, discounted
at the risk-free rate plus the hazard frozen at its level at $H$. With $H=40$ years and
$r\approx5\%$, the tail is discounted by $e^{-rH}\approx0.13$ and weighted by the
probability of survival without call to year 40, which is below $3\%$ in the base case
(Table~\ref{tab:callext}); its contribution to the price is a few hundredths of a point.

\subsection{Conditional Monte Carlo}

The write-down time enters \eqref{eq:price} only through the indicators
$\1_{\{t<\tau_P\}}$. Every other quantity is a functional of the state paths.

\begin{theorem}[Conditional pricing formula]\label{thm:cmc}
Let $\mathrm{CF}_i$ be the coupon cash flows of \eqref{eq:price} with the indicator of
$\tau_P$ removed, and let $\mathrm{CF}_{\mathrm{end}}$ and $T_{\mathrm{end}}$ be the
redemption, conversion or tail payment and its date on each path. Then
\begin{equation}\label{eq:cmc}
  V_0(\lambda)=\E\Biggl[\,\sum_{i}\mathrm{CF}_i\,e^{-rt_i-\lambda I_{t_i}}
  + \mathrm{CF}_{\mathrm{end}}\,e^{-rT_{\mathrm{end}}-\lambda I_{T_{\mathrm{end}}}}\Biggr].
\end{equation}
\end{theorem}

\begin{proof}
All of $\mathrm{CF}_i$, $\mathrm{CF}_{\mathrm{end}}$, $T_{\mathrm{end}}$, $\tau_m$ and
$\tau_c$ are $\G_\infty$-measurable. By the tower property,
$\E[X\1_{\{t<\tau_P\}}]=\E[X\,\Q(\tau_P>t\mid\G_\infty)]=\E[Xe^{-\lambda I_t}]$ for any
$\G_\infty$-measurable integrable $X$, by \eqref{eq:survat1}. Apply this to each term of
\eqref{eq:price}. In the tail term the event $\{\tau_P>H\}$ contributes $e^{-\lambda I_H}$,
and the value of the perpetuity after $H$ is given by \eqref{eq:tail}.
\end{proof}

The estimator replaces the simulated indicator by its conditional expectation. Two
consequences follow.

\begin{proposition}[Variance reduction]\label{prop:var}
Let $X$ be the discounted payoff of \eqref{eq:price} on one path, with $\tau_P$
simulated, and $Y=\E[X\mid\G_\infty]$ the integrand of \eqref{eq:cmc}. Then
$\E[Y]=\E[X]$ and $\operatorname{Var}(Y)=\operatorname{Var}(X)-\E[\operatorname{Var}(X\mid\G_\infty)]\le\operatorname{Var}(X)$,
with equality only if $X$ is $\G_\infty$-measurable.
\end{proposition}

\begin{proof}
This is the law of total variance.
\end{proof}

\begin{proposition}[Shape of the price in the hazard]\label{prop:shape}
On any fixed set of paths, and in expectation, the map $\lambda\mapsto V_0(\lambda)$ on
$[0,\infty)$ is infinitely differentiable, strictly decreasing and strictly convex, with
\begin{equation}\label{eq:dv}
  V_0'(\lambda) = -\E\Bigl[\sum_i\mathrm{CF}_i\,I_{t_i}e^{-rt_i-\lambda I_{t_i}}
  + \mathrm{CF}_{\mathrm{end}}\,I_{T_{\mathrm{end}}}e^{-rT_{\mathrm{end}}-\lambda I_{T_{\mathrm{end}}}}\Bigr]
  + \text{(tail term)} <0,
\end{equation}
$V_0(\lambda)\to0$ as $\lambda\to\infty$. Consequently, for any market price
$P\in(0,V_0(0))$ the implied hazard $\lambda^\star$ solving $V_0(\lambda^\star)=P$ exists
and is unique, and any bracketing method on $[0,\bar\lambda]$ with
$V_0(\bar\lambda)<P$ converges to it.
\end{proposition}

\begin{proof}
Each summand has the form $a\,e^{-\lambda b}$ with $a\ge0$ and $b=I_t>0$ for $t>0$. The
driver starts at $s_0/s_0=1$ (when CET1 is above the threshold) and is continuous, so
$I_t>0$ for every $t>0$ on every path. Each summand is therefore smooth, decreasing and
convex in $\lambda$, strictly so when $a>0$. The tail summand is the product of
$(y^{\mathrm{ref}}+m)F/(r+\lambda d)$ and $e^{-\lambda I_H}$ with $d=s_H^+/s_0\ge0$; both
factors are positive, decreasing and convex, and for two such functions
$(gh)''=g''h+2g'h'+gh''\ge0$ since $g'h'\ge0$. Sums and averages preserve these
properties, and at least one coupon is positive with positive probability, which gives
strictness. Every summand vanishes as $\lambda\to\infty$; dominated convergence gives the
limit. Existence and uniqueness of the root follow from continuity and strict
monotonicity.
\end{proof}

In practice we solve $V_0(\lambda)=P^{\mathrm{clean}}+\mathrm{AI}$ with Brent's method
\citep{Brent1973} on $[0,1]$. Because the path-wise cash flows and the integrated
drivers are stored once, each evaluation of $V_0(\lambda)$ is a vectorised sum, and the
root costs a few milliseconds. Proposition~\ref{prop:shape} also gives the price
sensitivity to the hazard, which we use in Section~\ref{sec:calibration} to translate
hazard differences into price units: for the three AT1s $|V_0'|$ is $0.077$--$0.084$
points per basis point of hazard.

\subsection{Monte Carlo error of the implied hazard}

Let $\hat V_n(\lambda)$ be the sample mean of $n$ paths. By the delta method, the
standard error of the implied hazard is approximately
\begin{equation}\label{eq:sehazard}
  \operatorname{se}(\hat\lambda^\star)\approx\frac{\operatorname{se}\bigl(\hat V_n(\lambda^\star)\bigr)}{\bigl|V_0'(\lambda^\star)\bigr|}.
\end{equation}
With $n=20{,}000$ paths, the standard error of the dirty price is $0.05$--$0.06$ points
and that of the implied hazard is $0.67$bp for all three bonds
(Appendix~\ref{app:numerics}). Because the three bonds are priced on the same paths, their
Monte Carlo errors are strongly positively correlated, and the cross-bond dispersion of
the implied hazards is more precise than each hazard: across independent seeds and path
counts from $5{,}000$ to $40{,}000$, the dispersion stays within $2.0$--$3.2$bp. The time
step matters more than the number of paths: on a fortnightly grid the dispersion is
$4.6$bp (Appendix~\ref{app:numerics}).

\subsection{Tier~2 and senior bonds}\label{sec:t2}

A Tier~2 or senior non-preferred note with coupons $c_i$ to a call date $T$ treated as
certain, recovery $R$ at resolution and resolution hazard $\lres s_t^+/s_0$ pays its
coupons and principal if $\tau_{\mathrm{res}}>T$, and $RF$ at $\tau_{\mathrm{res}}$ if
$\tau_{\mathrm{res}}\le T$.

\begin{proposition}[Tier~2 price]\label{prop:t2}
With recovery paid at the resolution time,
\begin{equation}\label{eq:t2cont}
  B_0=\E\Biggl[\sum_{t_i\le T}c_iF\,e^{-rt_i-\lres I_{t_i}}+F\,e^{-rT-\lres I_T}
  + RF\int_0^Te^{-ru}\,\lres\frac{s_u^+}{s_0}\,e^{-\lres I_u}\dd u\Biggr].
\end{equation}
On a grid $0=u_0<u_1<\dots<u_N=T$ with recovery paid at the end of the step in which
resolution occurs, the recovery leg is exactly
\begin{equation}\label{eq:t2grid}
  RF\sum_{n=0}^{N-1}\Bigl(e^{-\lres I_{u_n}}-e^{-\lres I_{u_{n+1}}}\Bigr)e^{-ru_{n+1}} ,
\end{equation}
and the difference from \eqref{eq:t2cont} is bounded by $RF\bigl(1-e^{-r\Delta}\bigr)$.
\end{proposition}

\begin{proof}
The survival terms follow as in Theorem~\ref{thm:cmc} with \eqref{eq:survt2}. For the
recovery leg, conditionally on $\G_\infty$ the law of $\tau_{\mathrm{res}}$ has density
$\lres(s_u^+/s_0)e^{-\lres I_u}$, so
$\E[RFe^{-r\tau_{\mathrm{res}}}\1_{\{\tau_{\mathrm{res}}\le T\}}\mid\G_\infty]$ is the
integral in \eqref{eq:t2cont}. On the grid, the conditional probability of resolution in
$(u_n,u_{n+1}]$ is $e^{-\lres I_{u_n}}-e^{-\lres I_{u_{n+1}}}$, which gives
\eqref{eq:t2grid}. Paying at $u_{n+1}$ instead of at the resolution time within the step
changes the discount factor by at most a factor $e^{-r\Delta}$, and the total probability
of resolution is at most one.
\end{proof}

With a monthly grid and $r\approx5\%$ the bound is $0.4\%$ of the recovered amount, that
is about $0.01$ points for a recovery of $20\%$ and a resolution probability of about
$13\%$ to the call. The price \eqref{eq:t2grid} is again smooth and strictly
decreasing in $\lres$ for $R<1$, and the implied $\lres$ is unique. Crucially, the Tier~2
is priced on the \emph{same} paths and the same driver $I_t$ as the AT1s, so that $\lres$
and $\lambda$ are directly comparable.

\subsection{A term structure for the resolution hazard}\label{sec:piecewise}

A single level $\lres$ is calibrated to one Tier~2 bond. With two bonds of different
maturities we can bootstrap a piecewise-constant level,
$\lres(t)=\lambda_1\1\{t\le T^\ast\}+\lambda_2\1\{t>T^\ast\}$, where $T^\ast$ is the call of
the shorter bond. The integrated hazard becomes
$\int_0^t\lres(u)\frac{s_u^+}{s_0}\dd u=\lambda_1I_{t\wedge T^\ast}+\lambda_2\bigl(I_t-I_{T^\ast}\bigr)^+$.

\begin{lemma}\label{lem:piecewise}
Because $I$ is non-decreasing, $I_{t\wedge T^\ast}=\min(I_t,I_{T^\ast})$ and
$(I_t-I_{T^\ast})^+=\max(I_t-I_{T^\ast},0)$. Hence every price under the piecewise hazard
is computed from the stored values $I_{t_i}$ together with the single path-wise value
$I_{T^\ast}$.
\end{lemma}

The bootstrap is sequential: $\lambda_1$ is implied from the shorter bond, whose cash
flows end at $T^\ast$; $\lambda_2$ is then implied from the longer bond with $\lambda_1$
fixed; finally the AT1-specific level is implied from each AT1 with the piecewise
resolution hazard in place. Each step is a one-dimensional root with the monotonicity of
Proposition~\ref{prop:shape}.

\subsection{Discretisation}

The state is simulated on a monthly grid ($\Delta=1/12$) over $H=40$ years with
$20{,}000$ paths. The spread uses its exact Gaussian transition,
\begin{equation}
  S_{n+1}=\bar s+(S_n-\bar s)e^{-a\Delta}+\eta\sqrt{\frac{1-e^{-2a\Delta}}{2a}}
  \Bigl(\rho Z^C_n+\sqrt{1-\rho^2}\,Z^\perp_n\Bigr),
\end{equation}
where $Z^C_n$ is the normal of the CET1 step and $Z^\perp_n$ an independent normal. The
CET1 ratio uses an Euler step for the diffusion with a Bernoulli jump of probability
$1-e^{-\lambda_J\Delta}$ and lognormal size. The integrated driver uses the left-point rule,
$I_{n+1}=I_n+\Delta\,s^+_n/s_0$.

Three discretisation effects deserve comment. First, the Euler scheme for a linear SDE has
weak order one; with $\kappa\Delta\approx0.03$ the bias in the transition mean is of order
$(\kappa\Delta)^2/2\approx5\times10^{-4}$ in relative terms. Second, the trigger is
monitored only at grid dates, which underestimates the probability of crossing between
dates. \citet{BroadieGlassermanKou1997} show that, for a Brownian motion with volatility
$\sigma$, discrete monitoring is approximately equivalent to shifting the barrier away by
$0.5826\,\sigma\sqrt{\Delta}$; here the shift is $0.08$pp of CET1 against a distance of
$7.8$pp to the trigger, and the conversion component of the spread is only $4$bp, so the
effect is negligible. Third, events on the grid are read at grid dates, while coupons are
discounted at their exact dates. Halving the step to fortnightly changes the implied
hazards by at most $1.6$bp, but not by the same amount for every bond, so the cross-bond
dispersion rises from $2.6$ to $4.6$bp (Appendix~\ref{app:numerics}).

\subsection{Common random numbers}

All primitive random numbers are drawn once and reused across bonds, scenarios, bumps
and feature switches. A difference between two prices computed on the same paths
(a Greek, a stress P\&L, a component of the spread decomposition, the dispersion of the
implied hazards across bonds) therefore carries no Monte Carlo noise of its own beyond
the second-order effect of the change on the path-wise payoffs. This is what makes the
attribution of Section~\ref{sec:attribution} exact on a given set of paths.